\subsection{向量值函数的积分的定义}

回顾一下，称 X 到 E 中的一个映射 \(\mathbf{f}\) 是弱连续的，如果对每个 \(\mathbf{z}' \in \mathrm{E}'\) ，映射 \(x \mapsto \langle \mathbf{f}(x), \mathbf{z}' \rangle\) （X 到 \(\mathbf{C}\) 中的；换言之即映射 \(\mathbf{z}' \circ \mathbf{f}\) ，也记作 \(\langle \mathbf{f}, \mathbf{z}' \rangle\) ）是连续的。我们称 X 到 E 中的映射 \(\mathbf{f}\) 是标量意义具有紧支集的，如果对每个 \(\mathbf{z}' \in \mathrm{E}'\) ，映射 \(x \mapsto \langle \mathbf{f}(x), \mathbf{z}' \rangle\) 具有紧支集。我们以 \(\widetilde{\mathcal{K}}(\mathrm{X};\mathrm{E})\) 表示 X 到 E 中的弱连续且标量意义具有紧支集的映射所成的空间；显然 \(\widetilde{\mathcal{K}}(\mathrm{X};\mathrm{E}) \supset \mathcal{K}(\mathrm{X};\mathrm{E})\) ，但这两个空间未必相同（见下面的例 2）；然而当 E 是有限维时它们相等。

注意，在弱连续（相应地，标量意义具有紧支集）函数的定义中，E 的拓扑仅通过对偶 \(E'\) 的中介而起作用；因此，当把 E 的拓扑换成任何使对偶相同的局部凸拓扑时，这些函数所成之集并不改变。

若 E 与 F 是一对对偶的向量空间，则应注意：说 X 到 E 中的映射对 \(\sigma(\mathrm{E},\mathrm{F})\) 连续与说它是弱连续的，二者含义相同。

设 \(\mathbf{f}\) 是 \(X\) 到 \(E\) 中的一个弱连续且标量意义具有紧支集的映射，并设 \(\mu\) 是 \(X\) 上的一个测度；对每个 \(\mathbf{z}' \in E'\) ，我们有 \(\mathbf{z}' \circ \mathbf{f} \in \mathcal{K}(X)\) ；令

\[
\varphi (\mathbf {z} ^ {\prime}) = \int \langle \mathbf {f} (x), \mathbf {z} ^ {\prime} \rangle d \mu (x) = \mu (\mathbf {z} ^ {\prime} \circ \mathbf {f}).
\]

显然 \(\varphi\) 是 \(\mathbf{E}'\) 上的一个线性形式，因而是 \(\mathbf{E}^{\prime*}\) 的一个元素。

定义 1. --- 对每个函数 \(\mathbf{f} \in \widetilde{\mathcal{K}}(\mathrm{X};\mathrm{E})\) ，我们把 \(\mathbf{f}\) 关于 \(\mu\) 的积分（记作 \(\int \mathbf{f} d\mu\) ，或 \(\int \mathbf{f}(x) d\mu(x)\) ，或 \(\int \mathbf{f}\mu\) ，或 \(\int \mathbf{f}(x)\mu(x)\) ）理解为 \(\mathrm{E}'^*\) 中由下式定义的元素：

\[
\left\langle \int \mathbf {f} d \mu , \mathbf {z} ^ {\prime} \right\rangle = \int \left\langle \mathbf {f}, \mathbf {z} ^ {\prime} \right\rangle d \mu \quad \text {for all} \mathbf {z} ^ {\prime} \in \mathrm{E} ^ {\prime}.\tag{1}
\]

注意，即使 E 是 Hausdorff 的且 \(\mathbf{f} \in \mathcal{K}(\mathrm{X};\mathrm{E})\) ，也未必有 \(\int f d\mu \in E\) （习题 1；参见 No.~3）。

例子. --- 1) 设 E 是 C 上有限维的 Hausdorff 空间，于是若 \((\mathbf{e}_{i})_{1\leqslant i\leqslant n}\) 是 E 的一个基，则映射

\[
\left(\xi_ {1}, \dots \xi_ {n}\right) \mapsto \sum_ {i = 1} ^ {n} \xi_ {i} \mathbf {e} _ {i}
\]

是 \(\mathbf{C}^n\) 到 \(\mathbf{E}\) 上的一个同构。于是我们知道 \(\mathbf{E}\) 上每个线性形式都是连续的，换言之， \(\mathbf{E}'\) 与 \(\mathbf{E}^*\) （\(\mathbf{E}\) 的代数对偶）相同，且 \(\mathbf{E}'^*\) 可典范地等同于 \(\mathbf{E}\) 。设 \((\mathbf{e}_i')_{1 \leqslant i \leqslant n}\) 是 \(\mathbf{E}'\) 中与 \((\mathbf{e}_i)\) 对偶的基；为使映射 \(\mathbf{f}\) （\(\mathbf{X}\) 到 \(\mathbf{E}\) 中的）是弱连续且标量意义具有紧支集的，必须且只须诸函数 \(f_i = \mathbf{e}_i' \circ \mathbf{f}\) 是具紧支集的连续函数；于是有 \(\mathbf{f}(x) = \sum_{i=1}^{n} f_i(x) \mathbf{e}_i\) 对所有 \(x \in \mathbf{X}\) 成立，且

\[
\int \mathbf {f} d \mu = \sum_ {i = 1} ^ {n} \mu (f _ {i}) \mathbf {e} _ {i}.
\]

\begin{enumerate}
\def\labelenumi{\arabic{enumi})}
\setcounter{enumi}{1}
\tightlist
\item
  取 E 为 X 上测度所成的空间 \(\mathcal{M}(\mathrm{X};\mathbf{C})\) ，赋以淡拓扑（ \(\S1\) , No.~9）；于是 E 的对偶 \(E'\) 可典范地等同于空间 \(\mathcal{K}(\mathrm{X};\mathbf{C})\) （TVS, II, \(\S6\) , No.~2, 命题 3）。X 到 E 中的映射 \(x \mapsto \varepsilon_{x}\) 是连续的（ \(\S1\) , No.~9, 命题 13），但当 X 不是紧时其支集不是紧的；然而它是标量意义具有紧支集的，因为对每个函数 \(f \in E'\) ，函数 \(x \mapsto \langle \varepsilon_x, f \rangle = f(x)\) 按定义具有紧支集。此外，
\end{enumerate}

\[
\int \left\langle \varepsilon_ {x}, f \right\rangle d \mu = \int f (x) d \mu (x) = \left\langle \mu , f \right\rangle
\]

对每个函数 \(f \in \mathcal{K}(\mathrm{X};\mathbf{C}) = \mathrm{E}'\) 成立，这就证明

\[
\int \varepsilon_ {x} d \mu (x) = \mu
\]

对 X 上的每个测度 \(\mu\) 成立。

\begin{enumerate}
\def\labelenumi{\arabic{enumi})}
\setcounter{enumi}{2}
\tightlist
\item
  若 \(\mathbf{E}\) 是 Hausdorff 的，则对每点 \(y \in \mathbf{X}\) 及每个函数 \(\mathbf{f} \in \widetilde{\mathcal{K}}(\mathbf{X};\mathbf{E})\) ，有
\end{enumerate}

\[
\int \mathbf {f} d \varepsilon_ {y} = \mathbf {f} (y)
\]

因为按定义 \(\int \langle \mathbf{f},\mathbf{z}^{\prime}\rangle d\varepsilon_{y} = \langle \mathbf{f}(y),\mathbf{z}^{\prime}\rangle\) 。

注. --- 1) 设 E 是一个局部凸空间，N 是 \(\{0\}\) 在 E 中的闭包，于是 \(E_{1} = E/N\) 是与 E 相伴的 Hausdorff 局部凸空间；我们知道对偶 \(E'\) 与 \(E_{1}'\) 相同；为使函数 f 属于 \(\widetilde{\mathcal{K}}(X;E)\) ，必须且只须 \(f_{1} = \pi \circ f\) （其中 \(\pi : E \to E_{1}\) 是典范同态）属于 \(\widetilde{\mathcal{K}}(X;E_{1})\) ，且此时 \(\int f d\mu = \int f_{1} d\mu\) 。因此可以只限于考虑 Hausdorff 局部凸空间。

\begin{enumerate}
\def\labelenumi{\arabic{enumi})}
\setcounter{enumi}{1}
\tightlist
\item
  设 E 是 C 上的一个局部凸空间， \(E_{0}\) 是 E 的底层 R 上局部凸空间；我们知道映射 \(z^{\prime} \mapsto Rz^{\prime}\) ——它把 E 上每个连续（复）线性形式 \(z^{\prime}\) 对应到连续（实）线性形式 \(z \mapsto R\langle z, z^{\prime} \rangle\) （\(E_{0}\) 上的）——是 \(E^{\prime}\) 到对偶 \(E_{0}^{\prime}\) （\(E_{0}\) 的）上的一个 R-同构（TVS, II, §8, No.~1）。类似地， \(E_{0}^{\prime *}\) （实向量空间 \(E_{0}^{\prime}\) 的代数对偶）可典范地等同于 \(E^{\prime *}\) （\(E^{\prime}\) 的代数对偶）的底层实空间。由此可知，若 \(\mu\) 是一个实测度而 f 是 \(\widetilde{\mathcal{K}}(X;E)\) 中的一个映射，则当把 f 看作取值于 \(E_{0}\) ，并把两个成员中出现的典范双线性型分别看作：第一个成员相对于 \(E_{0}^{\prime}\) 与 \(E_{0}^{\prime *}\) 之间的对偶，第二个成员相对于 \(E_{0}\) 与 \(E_{0}^{\prime}\) 之间的对偶时，公式 (1) 仍然成立。
\end{enumerate}

